\section{问题分析}
\subsection{切分、分核和顺序的耦合}
切图决定哪些操作共享一个驻留域。切得更细可能增加并行机会，却会让同一张量进入更多驻留域；合并操作可以减少结构搬运，却会增加同一核心的容量压力。即使切图和分核不变，核内顺序也会改变张量的存活区间，从而改变换入换出数量和完成时间。因此，本文不把三类决策拆成互不影响的独立优化，而是在每个候选方案上重新计算域、顺序和容量状态。

\subsection{问题一的难点}
场景 A 中每个子图是独立 Task，同核相邻 Task 之间也不能直接复用数据。任务激活同时受子图依赖、跨核等待和同核切换等待约束。单纯按计算周期均匀切分不能反映重复输入、跨 Task 写回和 DDR 争用；因此需要先用驻留域集合计算结构搬运，再用官方模拟确认并行与等待的实际平衡。

\subsection{问题二的难点}
场景 B 允许同核复用，但驻留数据会占用 L1/UB。一个候选方案可能结构搬运较少，却因为张量存活区间重叠而产生更多 Spill；也可能轻量估时略差，但通过改变末次消费者位置减少完整模拟中的额外搬运。本文将普通改善链和容量压力下降代表分开保存，让容量机制有独立的比较入口。

\subsection{问题三的难点}
共享 Cache 的效果由容量、带宽和查询时序共同决定。把同一输入的消费者全部聚合到一核可能减少重复查询，却可能损失计算并行；把消费者分散到多个核可能提高并行度，却会让多次查询在前一次填充完成前同时 miss。本文在同一问题二方案上分别构造聚合、分散和时序候选，以 Cache 事件回放的结果决定是否进入官方完整评估。

\subsection{方法选择}
候选搜索采用“结构可解释的起点 + 有界局部编辑 + 分层评估”。四类起点分别对应整图、连通分量装箱、关键路径切块和优先释放容量的切块；局部编辑包括迁移、拆分、合并、边界操作调整和同核调序。轻量层只用于筛选，展开层负责核内顺序和容量，完整层负责官方 Makespan。这样可以把每个改进对应到一个可检查的机制，同时控制搜索预算。
